% !TeX program = pdflatex
% !TeX encoding = UTF-8
% !BIB program = bibtex
% !TeX root = main.tex
\documentclass[conference]{IEEEtran}
\IEEEoverridecommandlockouts

\usepackage{cite}
\usepackage{amsmath,amssymb,amsfonts}
\usepackage{algorithm}
\usepackage{algorithmic}
\usepackage{graphicx}
\usepackage{textcomp}
\usepackage{xcolor}
\usepackage{booktabs}
\usepackage{multirow}
\usepackage{subfig}
\def\BibTeX{{\rm B\kern-.05em{\sc i\kern-.025em b}\kern-.08em
    T\kern-.1667em\lower.7ex\hbox{E}\kern-.125emX}}

\graphicspath{{figures/}}

\begin{document}

\title{Integrated Sensing and Communication for Multi-UAV Systems:
A Placeholder Title\\
% Sub-title should not be used for IEEE Xplore
\thanks{Identify applicable funding agency here. If none, delete this.}
}

\author{\IEEEauthorblockN{Dankao Chen}
\IEEEauthorblockA{\textit{Department of ...} \\
\textit{University of ...}\\
City, Country \\
email@example.com}
\and
\IEEEauthorblockN{2\textsuperscript{nd} Given Name Surname}
\IEEEauthorblockA{\textit{Department of ...} \\
\textit{University of ...}\\
City, Country \\
email@example.com}
}

\maketitle

\begin{abstract}
This work investigates integrated sensing and communication (ISAC) in
multi-UAV (MUAV) systems. Replace this placeholder abstract with a concise
summary of the problem, the proposed approach, and the key results. Do not
use symbols, special characters, footnotes, or math in the abstract.
\end{abstract}

\begin{IEEEkeywords}
integrated sensing and communication (ISAC), multi-UAV systems, beamforming,
resource allocation.
\end{IEEEkeywords}

\section{Introduction}
Background and motivation of ISAC in multi-UAV systems. State the
contributions of this work, typically as a bullet list:
\begin{itemize}
\item We formulate the MUAV-assisted ISAC problem as ...
\item We propose a ... algorithm to ...
\item Simulation results demonstrate that ...
\end{itemize}

\subsection{Related Work}
Review prior work on ISAC, UAV communications, and multi-UAV cooperation,
and position this paper against existing contributions~\cite{ref_isac_survey,
ref_uav_com}.

\subsection{Contributions and Organization}
Summarize the contributions and outline the remainder of the paper.
Section~\ref{sec:system} describes the system model.
Section~\ref{sec:method} presents the proposed method.
Section~\ref{sec:simulation} reports simulation results, and
Section~\ref{sec:conclusion} concludes the paper.

\section{System Model}\label{sec:system}

\subsection{Network Topology}
Describe the scenario: UAV positions, users/targets, channels.
Fig.~\ref{fig:scenario} illustrates the considered system.

\begin{figure}[!t]
\centering
% \includegraphics[width=\linewidth]{scenario.pdf}
\framebox[0.9\linewidth]{\rule{0pt}{2.5cm} placeholder}
\caption{Illustration of the MUAV-assisted ISAC scenario.}
\label{fig:scenario}
\end{figure}

\subsection{Communication Model}
Define the communication channels and rate expressions, e.g.,
\begin{equation}\label{eq:rate}
R_{k} = \log_{2}\!\left(1 + \frac{\left|\mathbf{h}_{k}^{H}\mathbf{w}_{k}\right|^{2}}
{\sum_{j\neq k}\left|\mathbf{h}_{k}^{H}\mathbf{w}_{j}\right|^{2} + \sigma^{2}}\right).
\end{equation}

\subsection{Sensing Model}
Define the sensing (e.g., radar) model, beampattern gain, and CRB, e.g.,
\begin{equation}\label{eq:beampattern}
G(\theta) = \mathbf{a}^{H}(\theta)\,\mathbf{R}_{x}\,\mathbf{a}(\theta),
\end{equation}
where $\mathbf{a}(\theta)$ denotes the steering vector and
$\mathbf{R}_{x}$ the transmit covariance matrix.

\section{Problem Formulation}
State the optimization problem, e.g.,
\begin{subequations}\label{opt:main}
\begin{align}
\max_{\{\mathbf{w}_{k}\}} \quad & \sum_{k} R_{k} \label{opt:obj}\\
\text{s.t.}\quad
& \mathrm{tr}(\mathbf{R}_{x}) \le P_{\max}, \label{opt:power}\\
& \mathrm{CRB}(\theta) \le \Gamma, \label{opt:crb}\\
& \|\mathbf{q}_{m} - \mathbf{q}_{m'}\|^{2} \ge d_{\min}^{2},
\quad \forall m \neq m'. \label{opt:collision}
\end{align}
\end{subequations}

\section{Proposed Method}\label{sec:method}
Describe the approach to solve problem~\eqref{opt:main}. An overall
algorithm is summarized in Algorithm~\ref{alg:proposed}.

\begin{algorithm}[!t]
\caption{Proposed Algorithm for MUAV-ISAC}
\label{alg:proposed}
\begin{algorithmic}[1]
\STATE Initialize UAV positions $\{\mathbf{q}_{m}^{(0)}\}$ and beams
\REPEAT
\STATE Update beamforming vectors by solving~\eqref{opt:main} with fixed UAV positions
\STATE Update UAV trajectories based on the obtained beams
\UNTIL convergence in objective~\eqref{opt:obj}
\STATE \textbf{return} optimized beams and trajectories
\end{algorithmic}
\end{algorithm}

\section{Simulation Results}\label{sec:simulation}

\subsection{Simulation Setup}
Describe parameters: carrier frequency, UAV altitude, noise power, path-loss
model, and benchmarks.

\subsection{Numerical Results}
Present the results. Table~\ref{tab:comparison} gives an example of a
three-point IEEE-style table.

\begin{table}[!t]
\caption{An Example Table}
\label{tab:comparison}
\centering
\begin{tabular}{@{}lccc@{}}
\toprule
Scheme & Metric 1 & Metric 2 & Metric 3\\
\midrule
Benchmark A & -- & -- & --\\
Benchmark B & -- & -- & --\\
Proposed    & -- & -- & --\\
\bottomrule
\end{tabular}
\end{table}

\section{Conclusion}\label{sec:conclusion}
Conclude the paper and outline possible future work.

\section*{Acknowledgment}
The authors would like to thank...

\bibliographystyle{IEEEtran}
\bibliography{references}

\end{document}
